% this text added by Adam on May 13, 2022
% this could end up as final text, but it deserves discussion at a DR telecon
% I think not enough of the DR telecon discussion has made it into the paper, so I decided it would be better to write this up as if it were a writeup and we can decide to cut it later

\section{Special cases}

Several regions required additional special attention in early 2022 because \tclean failed for a variety of reasons related to software configurations.
We were never able to fully determine what the underlying causes of the failures were, but they manifested as \texttt{CASA} like the following:

\noindent 
\texttt{MPICommandServer:: 
command\_request\_handler
\_service:: SynthesisImagerVi2:: CubeMajorCycle::MPIServer-23 (file src/code/synthesis/ImagerObjects/ CubeMajorCycleAlgorithm.cc, line 336)      Exception for chan range [1921, 1921] ---   Error in running Major Cycle : A nasty Visbuffer2 error occured...wait for CNGI
W51-E\_B3\_fullcube\_12M\_0\_34232805.log: Exception: Error in running Major Cycle : A nasty Visbuffer2 error occured...wait for CNGI}

or

\noindent
\texttt{SEVERE  tclean::::casa  Task tclean raised an exception of class RuntimeError with the following message: Error in making PSF : One or more  of the cube section failed in de/gridding. Return values for the sections: [1, 1, 1, 1, 1, 1, 1, 0, 0, 1, 1, 1, 1, 1, 1, 1, 1, 1, 1, 1, 1, 1, 1, 1, 1, 1, 1, 1, 1, 1, 1, 1, 1, 1, 1, 1, 1, 1, 1, 1, 1, 1, 1, 1, 1, 1, 1, 1, 1, 1, 1, 1, 1, 1, 1, 1, 1, 1, 1, 1, 1, 1, 1, 1, 1, 1, 1, 1, 1, 1, 1, 1, 1, 1, 1, 1, 1, 1, 1, 1, 1, 1, 1, 1, 1, 1, 1, 1, 1, 1, 1, 1, 1, 0, 1, 1, 1, 1, 1, 1, 1, 1, 1, 1, 1, 1, 1, 1, 1, 1, 1, 1, 1, 1, 1, 1, 1, 1, 1, 1, 1, 1, 1, 1]}
\newline
Two combinations of field and band, W51-E B3 spw0 and G333.60 B6 spw 1, proved particularly difficult to run \texttt{tclean} to completion.
We therefore modified the W51-E B3 spw0 imaging parameters to specify the number of channels by hand.
For G333.60, we found that the final model exhibited `zigzag' spectral patterns, in which, at the end of cleaning, adjacent spectral channels oscillated from $-1t$ to $+1t$, where $t$ is the final clean threshold.
These oscillations manifest strongly in the image and spectra and render the data unusable for science, even though technically the clean threshold has been met.
To correct this issue, we re-processed the final \texttt{.model} file by applying a 3-pixel maximum spectral filter, i.e., replacing the value at each pixel in the cube with the maximum of the $\pm1$ pixel range surrounding it in the spectral dimension.
This approach has the effect of retaining the spectral peaks, but slightly over-predicting the values to either edge of the peaks.